\documentclass[10pt]{article}
\usepackage[T1]{fontenc}
\usepackage[utf8]{inputenc}
\usepackage[french]{babel}
\usepackage{amsmath,amssymb,mathrsfs,mathpazo}
\usepackage{geometry}
\geometry{a4paper,margin=2cm}
\usepackage{graphicx,array,tabularx,tabularray,colortbl,multirow,multicol,booktabs,enumitem,pifont,ulem,fancybox,diagbox,cancel}
\usepackage{pstricks,pstricks-add,pst-plot,pst-tree,pst-node,pst-text,pst-func,pst-math,pst-bspline,pst-3dplot,pst-eucl,pst-solides3d}
\let\pscalculate\undefined
\usepackage{pst-bezier}
\usepackage{tikz,pgfplots,tkz-tab}
\usetikzlibrary{arrows,arrows.meta,patterns,calc,positioning,trees,decorations.markings,shapes.geometric,shapes.misc,shapes.arrows,intersections,decorations.pathmorphing,decorations.pathreplacing,backgrounds,shadows,babel}
\pgfplotsset{compat=1.18}
\usepgfplotslibrary{fillbetween}
\usepackage[np]{numprint}
\usepackage{listings,makecell,pas-tableur,twoopt}
\let\barre\undefined
\usepackage[pstricks]{tablvar}
\definecolor{darkgreen}{RGB}{0,100,0}
\providecommand{\red}{\color{red}}
\providecommand{\blue}{\color{blue}}
\providecommand{\green}{\color{green}}
\providecommand{\black}{\color{black}}
\providecommand{\white}{\color{white}}
\providecommand{\gray}{\color{gray}}
\providecommand{\cyan}{\color{cyan}}
\providecommand{\magenta}{\color{magenta}}
\providecommand{\yellow}{\color{yellow}}
\providecommand{\violet}{\color{violet}}
\providecommand{\orange}{\color{orange}}
\providecommand{\euro}{\texteuro}
\providecommand{\EUR}{\texteuro}
\providecommand{\ieme}{\textsuperscript{e}}
\providecommand{\degres}{\ensuremath{^\circ}}
\providecommand{\degree}{\ensuremath{^\circ}}
\providecommand{\vv}[1]{\overrightarrow{#1}}
\providecommand{\Coord}[1]{#1}
\providecommand{\decosix}{}
\providecommand{\starredbullet}{\ensuremath{\bullet}}
\pagestyle{empty}
\setlength{\parindent}{0pt}
\usepackage{adjustbox}\DeclareUnicodeCharacter{2236}{:}\DeclareUnicodeCharacter{2190}{\ensuremath{\leftarrow}}
\providecommand{\R}{\mathbb{R}}
\providecommand{\N}{\mathbb{N}}
\providecommand{\Z}{\mathbb{Z}}
\providecommand{\Q}{\mathbb{Q}}
\providecommand{\D}{\mathbb{D}}
\providecommand{\C}{\mathbb{C}}
\providecommand{\e}{{\mathrm{e}}}
\providecommand{\vect}[1]{\overrightarrow{#1}}
\providecommand{\vectt}[1]{\overrightarrow{\mathrm{#1}}}
\providecommand{\vvt}[1]{\overrightarrow{\mathrm{#1}}}
\providecommand{\np}[2][]{{#2~#1}}
\providecommand{\vv}[1]{{\overrightarrow{#1}}}
\renewcommand{\d}{{\mathrm{d}}}
\providecommand{\degres}{{^\circ}}
\providecommand{\degre}{{^\circ}}
\providecommand{\blue}{\color{blue}}
\providecommand{\red}{\color{red}}
\providecommand{\enskip}{\quad}
\providecommand{\llbracket}{\lbrack\!\lbrack}
\providecommand{\rrbracket}{\rbrack\!\rbrack}
\providecommand{\euro}{{\text{€}}}
\providecommand{\no}{{\text{n°}}}
\providecommand{\interval}[2]{{[#1;#2]}}
\providecommand{\intervFO}[2]{{[#1;#2[}}
\providecommand{\dotfill}{\ldots}

\providecommand{\ds}{\displaystyle}
\providecommand{\textding}[1]{#1}
\providecommand{\Oij}{\(\mathrm{O};\vec{\imath},\vec{\jmath}\)}
\providecommand{\Oijk}{\(\mathrm{O};\vec{\imath},\vec{\jmath},\vec{k}\)}
\providecommand{\Ouv}{\(\mathrm{O};\vec{u},\vec{v}\)}
\providecommand{\intervFF}[2]{\ensuremath{[#1;#2]}}
\renewcommand{\euro}{\ifmmode\text{\texteuro}\else\texteuro\fi}
\DeclareUnicodeCharacter{00B0}{\ifmmode^{\circ}\else\textdegree\fi}
\DeclareUnicodeCharacter{20AC}{\ifmmode\text{\texteuro}\else\texteuro\fi}
% Énoncés : Éducation nationale. Transcription et corrigés : APMEP, Denis Vergès et contributeurs. Découpage et classement : CoopMaths, Sébastien Lozano, Éric Elter et contributeurs. Adaptation de la présentation : Maths Explorer.
% CC BY-SA : https://coopmaths.fr/www/about/
% https://coopmaths.fr/alea/static/bac/2023/tex/bac_2023_03_sujet2_metropole_4.tex
\begin{document}
\noindent\textit{APMEP / CoopMaths. CC BY-SA. Présentation adaptée par Maths Explorer.}\par\medskip
\medskip

On considère la fonction $f$ définie sur $\R$ par 
\[f(x) = \ln \left(1 + \text{e}^{-x}\right),\]

où ln désigne la fonction logarithme népérien.
 
On note $\mathcal{C}$ sa courbe représentative dans un repère orthonormé \Oij. 

La courbe $\mathcal{C}$ est tracée ci-dessous.

\begin{center}
\psset{unit=2.25cm,yunit=2.25cm,arrowsize=2pt 3}
\begin{pspicture*}(-3,-0.8)(2.4,2.6)
\psaxes[linewidth=1.25pt,Dx=10,Dy=10]{->}(0,0)(-3,-0.8)(2.4,2.6)
\psaxes[linewidth=1.25pt](0,0)(1,1)
\psplot[plotpoints=2000,linecolor=red,linewidth=1.25pt]{-3}{2.4}{2.71828 x neg exp 1 add ln}
\psplotTangent[linestyle=dashed]{0}{4}{2.71828 x neg exp 1 add ln}
\psline[linestyle=dashed](-3,2.7)(3.4,-0.45)
\psdots(-2.2,2.3)(2.3,0.1)
\uput[u](-2.2,2.3){$M_a$}\uput[u](2.3,0.1){$N_a$}
\uput[dl](-1.6,1.8){\red $\mathcal{C}$}\uput[d](2.3,-0.43){$T_0$}
\end{pspicture*}
\end{center}

\medskip

\begin{enumerate}
\item 
	\begin{enumerate}
		\item Déterminer la limite de la fonction $f$ en $-\infty$.
		\item Déterminer la limite de la fonction $f$ en $+ \infty$.
		
Interpréter graphiquement ce résultat.
		\item On admet que la fonction $f$ est dérivable sur $\R$ et on note $f'$ sa fonction dérivée.
		
Calculer $f'(x)$ puis montrer que, pour tout nombre réel $x$,\: $f'(x) = \dfrac{-1}{1 + \text{e}^x}$.
		\item Dresser le tableau de variations complet de la fonction $f$ sur $\R$.
	\end{enumerate}	
\item On note $T_0$ la tangente à la courbe $\mathcal{C}$ en son point d'abscisse $0$.
	\begin{enumerate}
		\item Déterminer une équation de la tangente $T_0$.
		\item Montrer que la fonction $f$ est convexe sur $\R$.
		\item En déduire que, pour tout nombre réel $x$, on a : 
		\[f(x) \geqslant  - \dfrac12x + \ln (2).\]
		\end{enumerate}	
\item Pour tout nombre réel $a$ différent de 0, on note $M_a$ et $N_a$ les points de la courbe $\mathcal{C}$ d'abscisses respectives $-a$ et $a$. 

On a donc : $M_a\left(-a~;~ f(-a)\right)$ et $N_a\left(a~;~ f(a)\right)$.
	\begin{enumerate}
		\item Montrer que, pour tout nombre réel $x$, on a : $f(x) - f(-x) = - x$.
		\item En déduire que les droites $T_0$ et $\left(M_aN_a\right)$ sont parallèles.
	\end{enumerate}
\end{enumerate}

\end{document}
